% from ega4-2-fr.tex
\oldpage[IV-2]{137}

on conclut donc de (0, 16.3.4) que l'on a
\[
\dim_B(M)\leq\dim_{B'}(M')+1
\leq\bigl(\dim(A')+
\dim_{B\otimes_Ak}(M\otimes_Ak)\bigr)+1
=\dim(A)+\dim_{B\otimes_Ak}(M\otimes_Ak).
\]

Comme par hypothèse les membres extrêmes de ces inégalités sont égaux, on
a nécessairement: (i) $\dim_{B'}(M')=\dim_B(M)-1$, et comme $M$ est par
hypothèse un $B$-module de Cohen-Macaulay, cela signifie que $x$ est
$M$-régulier (0, 16.1.9 et 16.5.5); (ii)
$\dim_{B'}(M')=\dim(A')+
\dim_{B'\otimes_{A'}k}(M'\otimes_{A'}k)$; comme $M'$ est un $B'$-module de
Cohen-Macaulay en vertu de (i) et de (0, 16.1.9 et 16.5.5) et que $A'$ est
régulier, l'hypothèse de récurrence prouve que $M'$ est un $A'$-module plat;
on déduit alors de $(0_{\mathrm{III}}, 10.2.7)$ que $M$ est un $A$-module
plat, puisque $x$ est $M$-régulier par (i).

\subsection{Platitude et dimension projective}

\begin{proposition}[6.2.1]
\label{IV.2.s6.2.1-fr}
\begin{enumerate}
\item[(i)] Soient $A$, $B$ deux anneaux, $\varphi:A\to B$ un homomorphisme tel
que $B$ soit un $A$-module plat. Alors, pour tout $A$-module $M$, on a
\begin{equation}
\dim.\operatorname{proj}_B(M\otimes_AB)
\leq\dim.\operatorname{proj}_A(M).
\tag{6.2.1.1}\label{IV.2.s6.2.1.1-fr}
\end{equation}
\item[(ii)] Supposons en outre que $A$ soit un anneau noethérien, $B$ un
$A$-module fidèlement plat et $M$ un $A$-module de type fini; alors
\begin{equation}
\dim.\operatorname{proj}_B(M\otimes_AB)
=\dim.\operatorname{proj}_A(M).
\tag{6.2.1.2}\label{IV.2.s6.2.1.2-fr}
\end{equation}
\end{enumerate}
\end{proposition}

(i) On peut se borner au cas où
$n=\dim.\operatorname{proj}_A(M)$ est finie; il existe donc une résolution
gauche
\[
0\to P_n\to P_{n-1}\to\cdots\to P_0\to M\to0
\]
de $M$ par des $A$-modules projectifs; comme $B$ est un $A$-module plat, la
suite
\[
0\to P_n\otimes_AB\to P_{n-1}\otimes_AB\to\cdots\to
P_0\otimes_AB\to M\otimes_AB\to0
\]
est exacte; d'autre part les $P_i\otimes_AB$ sont des $B$-modules projectifs;
d'où la conclusion.

(ii) Supposons que $\dim.\operatorname{proj}_B(M\otimes_AB)=m$, et considérons
une suite exacte
\[
0\to R\to P_{m-1}\to P_{m-2}\to\cdots\to P_0\to M\to0
\]
où les $P_i$ $(0\leq i\leq m-1)$ sont des $A$-modules projectifs de type fini;
comme $A$ est noethérien, $R$ est aussi un $A$-module de type fini. Comme $B$
est un $A$-module plat, on a aussi une suite exacte
\[
0\to R\otimes_AB\to P_{m-1}\otimes_AB\to\cdots\to
P_0\otimes_AB\to M\otimes_AB\to0
\]
et l'hypothèse sur $M\otimes_AB$ entraîne que $R\otimes_AB$ est un $B$-module
projectif (0, 17.2.1). Comme $B$ est un $A$-module fidèlement plat, on en
conclut que $R$ est un $A$-module projectif (Bourbaki, \emph{Alg. comm.},
chap.~I\textsuperscript{er}, §~3, n\textsuperscript{o}~6, prop.~12); donc
$\dim.\operatorname{proj}_A(M)\leq m$, ce qui achève la démonstration.

\begin{corollary}[6.2.2]
\label{IV.2.s6.2.2-fr}
Soient $f:X\to Y$ un morphisme plat de préschémas, $\mathscr E$ un
$\mathcal O_Y$-Module quasi-cohérent. Si
$\dim.\operatorname{proj}(\mathscr E)\leq n$ (0, 17.2.14), on a
$\dim.\operatorname{proj}(f^*(\mathscr E))\leq n$.
\end{corollary}

